% Options for packages loaded elsewhere
\PassOptionsToPackage{unicode}{hyperref}
\PassOptionsToPackage{hyphens}{url}
%
\documentclass[
  12pt,
  a4paper,
  DIV=11,
  numbers=noendperiod]{scrreprt}

\usepackage{amsmath,amssymb}
\usepackage{setspace}
\usepackage{iftex}
\ifPDFTeX
  \usepackage[T1]{fontenc}
  \usepackage[utf8]{inputenc}
  \usepackage{textcomp} % provide euro and other symbols
\else % if luatex or xetex
  \usepackage{unicode-math}
  \defaultfontfeatures{Scale=MatchLowercase}
  \defaultfontfeatures[\rmfamily]{Ligatures=TeX,Scale=1}
\fi
\usepackage{lmodern}
\ifPDFTeX\else  
    % xetex/luatex font selection
\fi
% Use upquote if available, for straight quotes in verbatim environments
\IfFileExists{upquote.sty}{\usepackage{upquote}}{}
\IfFileExists{microtype.sty}{% use microtype if available
  \usepackage[]{microtype}
  \UseMicrotypeSet[protrusion]{basicmath} % disable protrusion for tt fonts
}{}
\usepackage{xcolor}
\setlength{\emergencystretch}{3em} % prevent overfull lines
\setcounter{secnumdepth}{5}
% Make \paragraph and \subparagraph free-standing
\makeatletter
\ifx\paragraph\undefined\else
  \let\oldparagraph\paragraph
  \renewcommand{\paragraph}{
    \@ifstar
      \xxxParagraphStar
      \xxxParagraphNoStar
  }
  \newcommand{\xxxParagraphStar}[1]{\oldparagraph*{#1}\mbox{}}
  \newcommand{\xxxParagraphNoStar}[1]{\oldparagraph{#1}\mbox{}}
\fi
\ifx\subparagraph\undefined\else
  \let\oldsubparagraph\subparagraph
  \renewcommand{\subparagraph}{
    \@ifstar
      \xxxSubParagraphStar
      \xxxSubParagraphNoStar
  }
  \newcommand{\xxxSubParagraphStar}[1]{\oldsubparagraph*{#1}\mbox{}}
  \newcommand{\xxxSubParagraphNoStar}[1]{\oldsubparagraph{#1}\mbox{}}
\fi
\makeatother


\providecommand{\tightlist}{%
  \setlength{\itemsep}{0pt}\setlength{\parskip}{0pt}}\usepackage{longtable,booktabs,array}
\usepackage{calc} % for calculating minipage widths
% Correct order of tables after \paragraph or \subparagraph
\usepackage{etoolbox}
\makeatletter
\patchcmd\longtable{\par}{\if@noskipsec\mbox{}\fi\par}{}{}
\makeatother
% Allow footnotes in longtable head/foot
\IfFileExists{footnotehyper.sty}{\usepackage{footnotehyper}}{\usepackage{footnote}}
\makesavenoteenv{longtable}
\usepackage{graphicx}
\makeatletter
\newsavebox\pandoc@box
\newcommand*\pandocbounded[1]{% scales image to fit in text height/width
  \sbox\pandoc@box{#1}%
  \Gscale@div\@tempa{\textheight}{\dimexpr\ht\pandoc@box+\dp\pandoc@box\relax}%
  \Gscale@div\@tempb{\linewidth}{\wd\pandoc@box}%
  \ifdim\@tempb\p@<\@tempa\p@\let\@tempa\@tempb\fi% select the smaller of both
  \ifdim\@tempa\p@<\p@\scalebox{\@tempa}{\usebox\pandoc@box}%
  \else\usebox{\pandoc@box}%
  \fi%
}
% Set default figure placement to htbp
\def\fps@figure{htbp}
\makeatother

\usepackage{indentfirst}
\usepackage{float}
\floatplacement{figure}{H}
\IfFileExists{plex-otf.sty}{
  %% Full TeXlive
  \usepackage[
    % math,
    RM={Scale=0.94},SS={Scale=0.94},SScon={Scale=0.94},TT={Scale=MatchLowercase,FakeStretch=0.9},
    DefaultFeatures={Ligatures=Common}
  ]{plex-otf}
  % \usepackage[Scale=MatchUppercase]{juliamono}
}{
  %% TinyTeX
  \usepackage{libertine}
}
\KOMAoption{captions}{tableheading}
\makeatletter
\@ifpackageloaded{caption}{}{\usepackage{caption}}
\AtBeginDocument{%
\ifdefined\contentsname
  \renewcommand*\contentsname{Содержание}
\else
  \newcommand\contentsname{Содержание}
\fi
\ifdefined\listfigurename
  \renewcommand*\listfigurename{Список иллюстраций}
\else
  \newcommand\listfigurename{Список иллюстраций}
\fi
\ifdefined\listtablename
  \renewcommand*\listtablename{Список таблиц}
\else
  \newcommand\listtablename{Список таблиц}
\fi
\ifdefined\figurename
  \renewcommand*\figurename{Рисунок}
\else
  \newcommand\figurename{Рисунок}
\fi
\ifdefined\tablename
  \renewcommand*\tablename{Таблица}
\else
  \newcommand\tablename{Таблица}
\fi
}
\@ifpackageloaded{float}{}{\usepackage{float}}
\floatstyle{ruled}
\@ifundefined{c@chapter}{\newfloat{codelisting}{h}{lop}}{\newfloat{codelisting}{h}{lop}[chapter]}
\floatname{codelisting}{Список}
\newcommand*\listoflistings{\listof{codelisting}{Листинги}}
\makeatother
\makeatletter
\makeatother
\makeatletter
\@ifpackageloaded{caption}{}{\usepackage{caption}}
\@ifpackageloaded{subcaption}{}{\usepackage{subcaption}}
\makeatother

\ifLuaTeX
\usepackage[bidi=basic]{babel}
\else
\usepackage[bidi=default]{babel}
\fi
\babelprovide[main,import]{russian}
\babelprovide[import]{english}
% get rid of language-specific shorthands (see #6817):
\let\LanguageShortHands\languageshorthands
\def\languageshorthands#1{}
\usepackage[style=gost-numeric,backend=biber,langhook=extras,autolang=other*]{biblatex}
\addbibresource{bib/cite.bib}
\usepackage{csquotes}
\usepackage{bookmark}

\IfFileExists{xurl.sty}{\usepackage{xurl}}{} % add URL line breaks if available
\urlstyle{same} % disable monospaced font for URLs
\hypersetup{
  pdftitle={Отчет по выполнению лабораторной работы №2},
  pdfauthor={Пашутина Анна Алексеевна},
  pdflang={ru-RU},
  hidelinks,
  pdfcreator={LaTeX via pandoc}}


\title{Отчет по выполнению лабораторной работы №2}
\usepackage{etoolbox}
\makeatletter
\providecommand{\subtitle}[1]{% add subtitle to \maketitle
  \apptocmd{\@title}{\par {\large #1 \par}}{}{}
}
\makeatother
\subtitle{Первоначальная настройка git}
\author{Пашутина Анна Алексеевна}
\date{}

\begin{document}
\maketitle

\renewcommand*\contentsname{Содержание}
{
\setcounter{tocdepth}{1}
\tableofcontents
}
\listoffigures
\listoftables

\setstretch{1.5}
\chapter{Цель
работы}\label{ux446ux435ux43bux44c-ux440ux430ux431ux43eux442ux44b}

Изучить идеологию и применение средств контроля версий. Освоить умения
по работе с git.

\chapter{Задание}\label{ux437ux430ux434ux430ux43dux438ux435}

Создать базовую конфигурацию для работы с git. Создать ключ SSH. Создать
ключ PGP. Настроить подписи git. Зарегистрироваться на Github. Создать
локальный каталог для выполнения заданий по предмету. \# Теоретическое
введение

\#Выполнение лабораторной работы.

Для начала установим git. В моем случае установлен (рис. 1)

\begin{figure}

{\centering \pandocbounded{\includegraphics[keepaspectratio]{image/image1.png}}

}

\caption{Установка git}

\end{figure}%

Теперь установим gh. (рис. 2)

\begin{figure}

{\centering \pandocbounded{\includegraphics[keepaspectratio]{image/image2.png}}

}

\caption{Установка gh}

\end{figure}%

Далее зададим имя для владельца репозитория. В данном случае это мои
фамилия и имя. (рис. 3)

\begin{figure}

{\centering \pandocbounded{\includegraphics[keepaspectratio]{image/image3.png}}

}

\caption{Указываю имя владельца репозитория}

\end{figure}%

Теперь зададим почту. Я задала ту почту, на которую зарегистрирован мой
аккаунт Github. (рис. 4)

\begin{figure}

{\centering \pandocbounded{\includegraphics[keepaspectratio]{image/image4.png}}

}

\caption{Указываю почту}

\end{figure}%

Настроим кодировку utf8 в выводе соообщений git. (рис. 5)

\begin{figure}

{\centering \pandocbounded{\includegraphics[keepaspectratio]{image/image5.png}}

}

\caption{Настраиваю кодировку utf8}

\end{figure}%

Зададим имя начальной ветки, настроим параметры autocrlf и safecrlf.
(рис. 6)

\begin{figure}

{\centering \pandocbounded{\includegraphics[keepaspectratio]{image/image6.png}}

}

\caption{Задаю имя начальной ветки и настриваю 2 параметра}

\end{figure}%

Создадим ключ RSA размером 4096 бит. (рис. 7)

\begin{figure}

{\centering \pandocbounded{\includegraphics[keepaspectratio]{image/image7.png}}

}

\caption{Создание ключа RSA}

\end{figure}%

Теперь создадим ключ по алгоритму ed25519 (рис. 8)

\begin{figure}

{\centering \pandocbounded{\includegraphics[keepaspectratio]{image/image8.png}}

}

\caption{Создание ключа ed25519}

\end{figure}%

Теперь создадим ключ gpg. Выбираем из предложенныч вариантов первый тип
(RSA или RSA), размер задаем 4096 бит, срок действия выбираем 0. (рис.
9)

\begin{figure}

{\centering \pandocbounded{\includegraphics[keepaspectratio]{image/image9.png}}

}

\caption{Создание ключа pgp}

\end{figure}%

После вводим свои данные: имя и адрес эл.почты;псоле соглашаемся с
генрацией ключа (рис. 10)

\begin{figure}

{\centering \pandocbounded{\includegraphics[keepaspectratio]{image/image10.png}}

}

\caption{Создание ключа pgp}

\end{figure}%

Далее выводим список pgp ключей (рис. 11)

\begin{figure}

{\centering \pandocbounded{\includegraphics[keepaspectratio]{image/image11.png}}

}

\caption{Выводим список ключей pgp}

\end{figure}%

Копируем наш ключ в буфер обмена (рис. 12)

\begin{figure}

{\centering \pandocbounded{\includegraphics[keepaspectratio]{image/image12.png}}

}

\caption{Копирование pgp ключа}

\end{figure}%

Вставляем этот ключ на github, задаем ему имя. У меня это имя Fedora
(рис. 13)

\begin{figure}

{\centering \pandocbounded{\includegraphics[keepaspectratio]{image/image13.png}}

}

\caption{Вставляю ключ в Github}

\end{figure}%

Теперь произведем настройку автоматических подписей (рис. 14)

\begin{figure}

{\centering \pandocbounded{\includegraphics[keepaspectratio]{image/image14.png}}

}

\caption{Настраиваю автоматические подписи}

\end{figure}%

Теперь нам надо было авторизоваться на Github с помощью gh. Выибираю
сайт для авторизации(Github), выбираю протокол(SSH), публичный SSH ключ
(id\_rsa.pub), имя для ключа(Fedora). В качестве способа авторизации
выбираю браузер (рис. 15)

\begin{figure}

{\centering \pandocbounded{\includegraphics[keepaspectratio]{image/image15.png}}

}

\caption{Авторизация в gh}

\end{figure}%

Теперь создаем рабочую директорию курса(у меня уже создана) и переходим
в нее с помощью cd. (рис. 16)

\begin{figure}

{\centering \pandocbounded{\includegraphics[keepaspectratio]{image/image16.png}}

}

\caption{Перехожу в директорию курса}

\end{figure}%

Далее создаем репозиторий для лабораторных работ из шаблона (рис. 17)

\begin{figure}

{\centering \pandocbounded{\includegraphics[keepaspectratio]{image/image17.png}}

}

\caption{Создаю репозиторий для лабораторных работ}

\end{figure}%

Клонируем его на свой компьютер (рис. 18)

\begin{figure}

{\centering \pandocbounded{\includegraphics[keepaspectratio]{image/image18.png}}

}

\caption{Клонирование репозитория}

\end{figure}%

Перходим в него с помощью cd и удаляем ненужные файлы(package.json) и
создаем необходимые каталоги, записав в файл COURSE строку os-intro(наш
текущий курс), прописываем make, чтобы нужные нам каталоги создались
(рис. 19)

\begin{figure}

{\centering \pandocbounded{\includegraphics[keepaspectratio]{image/image19.png}}

}

\caption{Удаляю ненужные файлы и ввожу комнаду make}

\end{figure}%

Теперь добавляем нашу папку для отправки с помщью git add . (рис. 20)

\begin{figure}

{\centering \pandocbounded{\includegraphics[keepaspectratio]{image/image20.png}}

}

\caption{Добавляю папку для отправки на github}

\end{figure}%

Деалем коммит, в котором указываем, что мы сделали структуру курса (рис.
21)

\begin{figure}

{\centering \pandocbounded{\includegraphics[keepaspectratio]{image/image21.png}}

}

\caption{Ипсользую git commit}

\end{figure}%

Отправляю файлы на сервер Github c помощью команды push (рис. 22)

\begin{figure}

{\centering \pandocbounded{\includegraphics[keepaspectratio]{image/image22.png}}

}

\caption{Использую git push}

\end{figure}%

\chapter{Выводы}\label{ux432ux44bux432ux43eux434ux44b}

Была произведена настройка git, проведена его первоначальная настройка,
были созданы елючи для авторизации и подписи, а также создан репозиторий
курса из предложенного шаблона.

\chapter{Ответы на контрольные
вопросы}\label{ux43eux442ux432ux435ux442ux44b-ux43dux430-ux43aux43eux43dux442ux440ux43eux43bux44cux43dux44bux435-ux432ux43eux43fux440ux43eux441ux44b}

\begin{enumerate}
\def\labelenumi{\arabic{enumi}.}
\tightlist
\item
  Системы контроля версий -- это системы, в которых мы можем хранить
  свои проекты и выкладывать их обновления, контролируя релизы и каждые
  внесённые изменения. Эти системы нужны для работы над проектами, чтобы
  иметь возможность контролировать версии проектов и в случае командной
  работы контролировать изменения, внесённые всеми участниками. Также,
  VCS позволяют откатываться на более ранние версии
\item
  Хранилище -- репозиторий, в нём хранятся все файлы проекта и все его
  версии\\
  commit -- внесённые изменения в репозитории\\
  история -- это история изменений файлов проекта\\
  рабочая копия -- копия, сделанная из версии репозитория, с которой
  непосредственно работает сам разработчик
\item
  Централизованные системы контроля версий имеют один центральный
  репозиторий, с которым работают все разработчики. Примером является
  CVS, который является уже устаревшей системой.\\
  В децентрализованных системах же используется множество репозиториев
  одного проекта у каждого из разработчиков, при этом репозитории можно
  объединять брать из каждого только то, что нужно. Примером является
  знакомый нам Git
\item
  Создаётся репозиторий, и разрабатывается проект. При внесении
  изменений файлы отправляются на сервер
\item
  Разработчик клонирует репозиторий к себе на компьютер, и после
  внесения изменений выгружает их на сервер в качестве отдельной версии.
  После этого разработчики с более высокими правами могут, например,
  объединить его версию с текущей
\item
  Хранение файлов проекта, а также обеспечение командной работы, и
  контроль за версиями проекта
\item
  git clone -- клонирует проект с сервера на компьютер\\
  git add -- добавляет папку для выгрузки на сервер\\
  git commit -- фиксирует изменения репозитория\\
  git push -- выгружает изменения на сервер\\
  git pull -- получить изменения с сервера\\
  git rm -- удалить файл\\
  git status -- получить статус репозитория
\item
  С локальным: git commit -am \enquote{added files} -- создаёт коммит С
  удалённым: git push -- загрузить данные на удалённый сервер
\item
  Ветки -- это несколько независимых копий проекта, в каждой из которых
  ведётся разработка какой-то конкретной функции, при этом ветки
  существуют параллельно. Они нужны, когда нужно параллельно вести
  разработку нескольких функций, а в конце их можно объединить в одну
\item
  Игнорировать файлы можно, внося их в файл .gitignore. Игнорировать
  файлы нужно, когда их не нужно добавлять в репозиторий. Например, это
  могут быть файлы виртуального окружения (venv)
\end{enumerate}

\chapter*{Список
литературы}\label{ux441ux43fux438ux441ux43eux43a-ux43bux438ux442ux435ux440ux430ux442ux443ux440ux44b}
\addcontentsline{toc}{chapter}{Список литературы}

\printbibliography[heading=none]





\end{document}
